\needspace{0.25\textheight}
\section*{Аннотация}
\addcontentsline{toc}{section}{Аннотация}

Работа даёт полное строгое изложение конечномерной модели динамики завихренности
Навье--Стокса на алгебре Ли $\so \cong \R^3$: геометрия проективного замыкания
$\RP^3 \cong \mathrm{SO}(3)$, вывод уравнений Эйлера--Пуанкаре с диссипацией,
инварианты Казимира и гамильтониан, каскад редукции $P^3 \to 2D \to 1D \to 0D$
(теорема Ладыженской, линеаризация Коула--Хопфа, спектральный проектор МакВини)
и главный результат об отсутствии сингулярностей в модели. Все утверждения
сопровождаются полными доказательствами; помимо исходной программы изложены
расширения: прямое доказательство тождества Якоби скобки Ли--Пуансона,
теорема периодичности орбит, оптимальность константы $4$ в оценке дефекта
идемпотентности МакВини и условная теорема о тензоре деформации, уточняющая
границы переноса результатов на полное уравнение Навье--Стокса. Раздел~8
содержит численную верификацию всех тождеств и оценок (23 проверки,
интегратор RK4): инварианты сохраняются с точностью $5{,}7\cdot 10^{-9}$
при порядке сходимости $p\approx 4{,}96$, орбиты замкнуты с точностью
$5\cdot 10^{-15}$, оценка Грёнуолла со штрейном достигается точно.

\bigskip

\needspace{0.25\textheight}
\section{Введение и постановка задачи}\label{sec:intro}

Пусть вязкая несжимаемая жидкость занимает область $\Omega \subset \R^3$,
поле скорости $\bu(x,t)$ удовлетворяет системе Навье--Стокса, а завихренность
$\bwn = \nabla \times \bu$~--- её вихревой форме
\begin{equation}\label{eq:ns-vort}
  \partial_t \bwn + (\bu \cdot \nabla)\bwn
  = (\bwn \cdot \nabla)\bu + \nu \Delta \bwn,
  \qquad \nabla \cdot \bu = 0 .
\end{equation}
Центральный вопрос глобальной регулярности \eqref{eq:ns-vort}~--- может ли норма
$\|\bwn(\cdot,t)\|_{L^\infty}$ достигнуть бесконечности за конечное время
$T^{*}$ \cite{fefferman}. Ключевую роль играет расщепление нелинейности:
обозначим $S = \tfrac12(\nabla\bu + (\nabla\bu)^{\mathsf T})$ тензор деформации
(симметричная часть $\nabla\bu$). Член растяжения вихря допускает тогда
точное представление $(\bwn\cdot\nabla)\bu = S\bwn$ (лемма~\ref{lem:strain}):
вихрь не «вращается», а только деформируется симметричной частью поля.
Именно этот член способен увеличивать $\|\bwn\|$, и именно он делает
трёхмерную проблему открытой.

\emph{Программа настоящей работы} (модель $P^3 \to \so$) состоит в замене
бесконечномерной динамики \eqref{eq:ns-vort} конечномерной системой
Эйлера--Пуанкаре на $\so^{*} \cong \R^3$ для амплитудного вектора
вихревой трубки $\Om = (\Omega_1,\Omega_2,\Omega_3)$:
\begin{equation}\label{eq:epd-intro}
  \dot{\Om} = \Om \times (\bbI^{-1}\Om) - \nu\,\bbL\,\Om,
\end{equation}
где $\bbI = \diag(I_1,I_2,I_3)$, $I_k > 0$~--- тензор инерции трубки,
$\bbL = \bbL^{\mathsf T} \succeq 0$~--- оператор вязкого затухания
(структурный аналог $-\Delta$), а нелинейность
$\Om \mapsto \Om \times (\bbI^{-1}\Om)$~--- чисто кососимметричный
поворот фазы, ортогональный $\Om$ при всех $t$. Модель сохраняет
геометрию задачи: состояние живёт на компактном многообразии
$\RP^3 \cong \mathrm{SO}(3)$, а $\Om$~--- в коалгебре $\so^{*}$.

Содержание работы~--- строгий разбор модели от геометрии до границ
применимости. Перечислим основные результаты; каждый пункт снабжён
ссылкой на точную формулировку.

\begin{enumerate}[label=\textup{(\roman*)}, leftmargin=2.2em, itemsep=2pt]
  \item \textbf{Геометрия} (\S\ref{sec:geom}): полный вывод диффеоморфизма
        $\RP^3 \cong \mathrm{SO}(3)$ через кватернионы, включая явную
        формулу поворота ось--угол (теорема~\ref{th:rp3so3}); тождества
        алгебры $\so \cong \R^3$ с метрикой Киллинга
        $\langle u,v\rangle = -\tfrac12\Tr(\hat u \hat v) = u\cdot v$
        (леммы~\ref{lem:hatjacobi}, \ref{lem:killing}).
  \item \textbf{Вывод системы} (\S\ref{sec:ep}): вариационный вывод
        \eqref{eq:epd-intro} из редуцированного лагранжиана
        (лемма~\ref{lem:variation}, теорема~\ref{th:ep}); скобка
        Ли--Пуансона $\{F,G\} = -\langle\Om, \nabla F \times \nabla G\rangle$
        с \emph{полным прямым доказательством} тождества Якоби
        (предложение~\ref{prop:poisson}); глобальная корректность задачи
        Коши (теорема~\ref{th:wp}).
  \item \textbf{Инварианты} (\S\ref{sec:inv}): точные законы эволюции
        Казимира $C = \|\Om\|^2$ и гамильтониана
        $\mathcal H = \tfrac12\langle\Om,\bbI^{-1}\Om\rangle$
        (теоремы~\ref{th:casimir}, \ref{th:ham}); экспоненциальное
        затухание $C(t) \le C(0)e^{-2\nu\lambda_{\min}t}$
        (следствие~\ref{cor:decay}); периодичность всех неравновесных
        орбит при $\nu = 0$ (предложение~\ref{prop:periodic});
        классификация и устойчивость равновесий, включая эффект
        «теннисной ракетки» (предложение~\ref{prop:stability}).
  \item \textbf{Каскад редукции} (\S\ref{sec:cascade}):
        точная лемма о растяжении $(\bwn\cdot\nabla)\bu = S\bwn$
        (лемма~\ref{lem:strain}); теорема Ладыженской об энстрофии
        (теорема~\ref{th:ladyzh}); полная линеаризация Коула--Хопфа
        (теорема~\ref{th:colehopf}); спектральная теория проектора
        МакВини $\mathcal Q(\bbP) = 3\bbP^2 - 2\bbP^3$ с доказательством
        \emph{оптимальности константы} $4$ в оценке дефекта
        (теорема~\ref{th:mcweeny}, предложение~\ref{prop:mcweeny-iter}).
  \item \textbf{Отсутствие сингулярностей} (\S\ref{sec:noblowup}):
        главный результат~--- априорная оценка
        $\sup_{t \ge 0}\|\Om(t)\| \le \|\Om(0)\| < \infty$ при всех
        $\nu \ge 0$ (теорема~\ref{th:main}); явные осесимметричные
        решения с вязким замедлением прецессии (предложение~\ref{prop:axisym}).
  \item \textbf{Границы метода} (\S\ref{sec:critique}): точное структурное
        объяснение, почему модель не решает трёхмерную проблему
        Клэя (несферичность коалгебраических орбит группы
        диффеоморфизмов против сферических орбит $\so$), и условная
        теорема со штрейном $\dot\Om = \ldots + S(t)\Om - \nu\bbL\Om$,
        воспроизводящая структуру критериев BKM/CKN
        (теорема~\ref{th:strain}).
  \item \textbf{Верификация} (\S\ref{sec:numerics}): численная проверка
        всех тождеств и оценок: 23/23 тестов, включая машинно-точную
        замкнутость орбит и точное достижение границы Грёнуолла.
\end{enumerate}

Мы сознательно отделяем доказанное внутри модели от вопроса о переносе
на уравнение \eqref{eq:ns-vort}: модельная теорема~\ref{th:main} полна и
строга, однако её перенос на трёхмерное уравнение требует контроля
$\|S(t)\|$, который модель постулирует, а полное уравнение не даёт;
подробное обсуждение~--- в \S\ref{sec:critique}. Нумерация всех величин
символьна: параметры $I_1, I_2, I_3$, $\nu$, $\bbL$ нигде не фиксируются,
кроме численных экспериментов \S\ref{sec:numerics}.
